\documentclass[11pt]{article}
\usepackage{amsmath,amssymb,amsthm,mathtools,fullpage}
\newtheorem{theorem}{Theorem}[section]
\newtheorem{proposition}[theorem]{Proposition}
\newtheorem{lemma}[theorem]{Lemma}
\newtheorem{corollary}[theorem]{Corollary}
\theoremstyle{definition}
\newtheorem{definition}[theorem]{Definition}
\theoremstyle{remark}
\newtheorem{remark}[theorem]{Remark}
\newcommand{\Z}{\mathbb{Z}}
\newcommand{\bl}{\boldsymbol{\lambda}}
\newcommand{\bs}{\mathbf{s}}
\newcommand{\Nlt}{N^{<}}
\newcommand{\ket}[1]{\lvert #1 \rangle}
\DeclareMathOperator{\Ind}{\mathbf 1}
\title{Q63: the vanishing branch.\\
A closed form for form (ii) of $[e_i,R^{(\ell)}(t)]$,\\
and the proof that it is empty at level one}
\author{Clio}
\date{3 September 2026}

\begin{document}
\maketitle

\begin{abstract}
The 2 September paper proved that every nonzero entry of $[e_i,R^{(\ell)}(t)]$ is either
$\varepsilon q^{k}t^{b}(1+q^{j}t)$ with $j$ known in closed form (form (i)), or
$t^{h}(q^{-N_1}-q^{-N_2})$ (form (ii)), and left $N_1-N_2$ uncomputed --- ``nobody has
looked''. We look.

The census first. All $1604$ form-(ii) entries of the standing sweep have exactly two
contributing paths, in \emph{opposite} commutator terms, with equal ribbon heights. They
split into three structural cases, and the split is the theorem: (A) the ribbon and the
node sit on different runners; (B) same runner, two beads moved; (C) same runner, one
bead moved, the two paths lying in opposite terms.

Then the closed form, in three parts. Case (B) \textbf{never occurs}: $N_1=N_2$
identically, so the entry is zero. Case (A) gives $N_1-N_2=\pm1$, read off a
\emph{two-site indicator} comparing the node's content with an end of the ribbon,
which end being selected by the tie-break order of the two runners. Case (C) is governed
by the $\sigma/\tau$ vocabulary of Theorem \texttt{thm:j}:
\[
\boxed{\;N_1-N_2\;=\;-\varepsilon\big(\sigma(\bl;c,d)+\tau(\bl;c+e,d)\big),\qquad
\varepsilon=+1\iff c-1\in M_d(\bl)\;}
\]
--- the same two half-tie-terms, at the same two contents, with the leading $1$ of
Theorem \texttt{thm:j} absent and the sign reversed. So the brief's instruction to test
the surviving vocabulary before inventing a new one was right for exactly half the
branch, and I say below which half and why.

Everything is proved from two lemmas: a \emph{ribbon-shift lemma} computing how a
Chevalley exponent changes when an $e$-ribbon is inserted underneath it (this does cases
(A) and (B)), and the telescope-difference computation of 2 September rerun with the two
paths in opposite terms (this does case (C)).

\textbf{Gap 2 closes.} Form (ii) is \emph{empty at $\ell=1$}, proved, not observed: case
(A) needs two runners, case (B) cancels always, and case (C) has $\sigma=\tau=0$ at
$\ell=1$. This is a cancellation and not a vacuity --- at $\ell=1$ there are $1424$ live
opposite-term path pairs in the checked range, every one of which sums to zero. The
node \texttt{Q63-formii-empty-at-level-1} is promoted from \texttt{computed} to
\texttt{proved}, and with it the classification of $[e_i,R^{(\ell)}(t)]$ is complete.

Verification: $1604/1604$ in sample; $1564/1564$ out of sample over seven configurations,
including $\ell=4$, which exhibits $|N_1-N_2|=3$, a value outside every in-sample range
and inside the one now proved. Thirteen negative controls, each checked non-degenerate
before being scored, all fail; the floor (the constant $N_1-N_2\equiv1$) is
$1530/3168$.

\textbf{Scope, stated here and not in a closing remark.} Proposition \texttt{prop:heis}
of 31 August shows $R^{(\ell)}(-q^{-1})\neq B^{[e]}_{-1}$ for $\ell\ge2$. What is
classified below is therefore the commutator of a concretely defined graded ribbon
operator, and \emph{not} known to be a statement about the level-$\ell$ Heisenberg
action on Fock space. Nothing in this paper narrows that gap.
\end{abstract}

\section{What was open, and what closes}\label{sec:summary}

Theorem \texttt{thm:dich} of 2 September classified every nonzero entry
$\Phi_{\nu\bl}(t)$ of $[e_i,R^{(\ell)}(t)]$ into

\begin{itemize}
\item \textbf{form (i)}: $\varepsilon q^{k}t^{b}(1+q^{j}t)$, the two paths in the same
commutator term, with
$j=1+\varepsilon\big(\sigma(\bl;c,d)+\tau(\bl;c+e,d)\big)$ (Theorem \texttt{thm:j});
\item \textbf{form (ii)}: $t^{h}(q^{-N_1}-q^{-N_2})$, the two paths in opposite terms,
vanishing at $q=1$, with $N_1-N_2$ \emph{uncomputed}.
\end{itemize}

Gaps 1 and 2 of that paper were: compute $N_1-N_2$; and decide whether form (ii) is
empty at $\ell=1$, where a $350$-entry sweep had found none but Theorem \texttt{thm:dich}
permitted them. Both close here, and they close together, because the closed form is what
decides the emptiness.

\paragraph{Order of work, and an honest note about it.} The brief asked for the census
before any formula. The census was run (\S\ref{sec:census}) and is reported first. But I
must record that I did not arrive at the formulas by looking at the census: reading the
proof of Lemma \texttt{lem:paths} and of Theorem \texttt{thm:j} on the way in, I saw that
the same two computations answer the other branch, and I wrote down the three predictions
of \S\ref{sec:closed} \emph{before} running \texttt{census.py}. That is a stronger
position than fitting, not a weaker one --- the census then tests a pre-registered
prediction rather than confirming a description of itself --- but it is a different
position from the one the brief specified, and the reader should know which they are
being offered. The census still did work the prediction could not: it is what shows that
case (B) pairs are \emph{numerous} rather than absent (\S\ref{sec:census}), which is the
difference between a theorem and a vacuity.

\section{Conventions}\label{sec:conv}

Those of 31 August and 2 September, unchanged: rank $n=e\ge2$, level $\ell\ge1$;
$\bl=(\lambda^{(1)},\dots,\lambda^{(\ell)})$ with multicharge $\bs$; runner-$d$ shifted
Maya set $M_d=\{\lambda^{(d)}_j-j+s_d\}$, with the deep tail (all positions below the
stored profile occupied) and the far top (none occupied). A node of component $d$ at
shifted content $x$ is the bead move $x-1\leftrightarrow x$ on runner $d$; an $e$-ribbon
is the bead move $b\mapsto b+e$, of height $h=\#\big(M_d\cap(b,b+e)\big)$. The tie-break
is $\prec$ with the standard sign: $(x,d')$ is \emph{below} $(x,d)$ iff $d'>d$. Write
$d'\prec d$ for ``runner $d'$ is below runner $d$''.

Throughout,
\[
\chi_d(x)\;=\;[\text{addable at }(d,x)]-[\text{removable at }(d,x)].
\]

\begin{lemma}[The abacus form of $\chi$]\label{lem:chi}
$\chi_d(x)=\Ind_{M_d}(x-1)-\Ind_{M_d}(x)$.
\end{lemma}
\begin{proof}
Addable at $(d,x)$ means $x-1\in M_d,\ x\notin M_d$; removable means
$x-1\notin M_d,\ x\in M_d$. Writing $a=\Ind_{M_d}(x-1)$, $b=\Ind_{M_d}(x)$, the two
indicators are $a(1-b)$ and $b(1-a)$, whose difference is $a-b$.
\end{proof}

This one-line identity is what makes everything below finite and local: $\chi_d$ is a
difference of indicators, so it telescopes, and it is supported on the finitely many
positions where $M_d$ is not constant.

\begin{definition}\label{def:sigmatau}
For $x\in\Z$ and a runner $d$,
\[
\tau(\bl;x,d)=\sum_{d'\prec d}\chi_{d'}(x),\qquad
\sigma(\bl;x,d)=\sum_{d'\succ d}\chi_{d'}(x),\qquad
A_{d'}(x)=\sum_{k\ge1}\chi_{d'}(x-ke),
\]
so that Theorem \texttt{thm:tele} reads
$\Nlt(\bl,\gamma)=\sum_{d'}A_{d'}(x)+\tau(\bl;x,d)$ for a node $\gamma$ at $(d,x)$.
Both $\sigma$ and $\tau$ read only runners $\ne d$, so both are unchanged if $\bl$ is
replaced by any multipartition agreeing with it off runner $d$. At $\ell=1$ both are
empty sums, hence $0$.
\end{definition}

Recall from Definition \texttt{def:chev} that the $e_i$ weight is $q^{-\Nlt}$ with
$\Nlt$ read on the \emph{target} of the removal. A \emph{path} is a ribbon move together
with a node removal, in one of the two orders; the $e_iR$ order contributes with sign
$+1$ and the $Re_i$ order with sign $-1$. For a form-(ii) entry with one path in each
term we always write $N_1$ for the exponent of the $e_iR$ path and $N_2$ for that of the
$Re_i$ path, so that $\Phi=t^{h}(q^{-N_1}-q^{-N_2})$.

\section{The census}\label{sec:census}

Run over the nine triples $(e,\ell,\bs)$ of Remark \texttt{rem:verif-rigid}, all
multipartitions of size $\le4$, all $i$; bead depth $|\bl|+2e+6$. The same run
reproduces $1248$ form-(i) entries and $0$ unclassified entries, which is the check that
this session's pipeline computes the object of the previous two sessions.

\begin{proposition}[Form-(ii) census]\label{prop:census2}
Over all $1604$ form-(ii) entries: every entry has exactly two contributing paths; in
every entry the two paths lie in \emph{opposite} commutator terms ($1604/1604$); in
every entry the two ribbon heights are equal ($1604/1604$); and every entry falls into
exactly one of
\begin{center}
\begin{tabular}{llrr}
\hline
case & description & count & $N_1-N_2$ observed\\
\hline
(A) & ribbon and node on different runners & $1146$ & $\pm1$ only\\
(B) & one runner, two beads moved & $\mathbf{0}$ & ---\\
(C) & one runner, one bead moved, opposite terms & $458$ & $\pm1,\pm2$\\
\hline
\end{tabular}
\end{center}
In cases (A) and (B) the two paths use the \emph{same} ribbon and the \emph{same} node,
in the two orders. In case (C) they use different ribbons and different nodes: the
$c$-path and the $(c+e)$-path of Lemma \texttt{lem:paths}.
\end{proposition}

Three things in this table are worth naming before they are proved.

\paragraph{The opposite-term fact is forced, and is the mirror of the form-(i) one.}
A form-(ii) entry is a single $t$-power whose coefficient vanishes at $q=1$, i.e.\ a
difference of two equal-signed monomials appearing with opposite signs; the two terms of
a commutator enter with opposite signs. (For form (i) the 2 September paper ran the same
argument the other way.)

\paragraph{Case (B) is empty, and that is a claim about entries, not about pairs.}
An entry-level census cannot distinguish ``case (B) always cancels'' from ``case (B)
never happens''. \emph{Variation exercised:} enumerate path \emph{pairs} rather than
entries, keeping those whose sum is zero, and ask separately whether case-(B) pairs occur
and whether they cancel.

\begin{center}
\begin{tabular}{lrrr}
\hline
configuration & case-(B) pairs & nonzero & $N_1-N_2$ \\
\hline
$e{=}2,\ell{=}2$, $|\bl|\le4$ & $60$ & $0$ & $0$ throughout\\
$e{=}3,\ell{=}2$, $|\bl|\le4$ & $94$ & $0$ & $0$ throughout\\
$e{=}5,\ell{=}2$, $|\bl|\le3$ & $72$ & $0$ & $0$ throughout\\
$e{=}2,\ell{=}3$, $|\bl|\le3$ & $36$ & $0$ & $0$ throughout\\
$e{=}3,\ell{=}3$, $|\bl|\le3$ & $66$ & $0$ & $0$ throughout\\
$e{=}2,\ell{=}4$, $|\bl|\le3$ & $56$ & $0$ & $0$ throughout\\
\hline
\end{tabular}
\end{center}

So case-(B) pairs are abundant --- $384$ of them here --- and every one of them cancels.
The emptiness is a theorem about a live population.

\paragraph{The case-(C) values reach $\ell-1$.} At $\ell=2$, $|N_1-N_2|\le1$; at
$\ell=3$ the value $2$ appears ($24$ times); the out-of-sample $\ell=4$ run produces $3$.
Case (A) never exceeds $1$ at any level. That is the shape the theorem below predicts,
and it is the first sign that (A) and (C) are governed by different statistics.

\paragraph{Census question 2: which object does each exponent read from?} $N_1$ is read
on the target $\nu$ of the entry; $N_2$ is read on the intermediate
$\rho=\bl-\gamma$, which is neither $\bl$ nor $\nu$. Their \emph{difference}, by
Theorems~\ref{thm:A} and \ref{thm:C} below, is a function of $\bl$ and the net bead
displacement alone --- intrinsic, exactly as $j$ turned out to be.

\paragraph{Census question 3: are $N_1$ and $N_2$ separately meaningful?} Yes, in the
weak sense that each is a well-defined finite integer: \emph{variation exercised}, the
abacus depth raised by $6$; over $590$ tracked pairs neither $N_1$ nor $N_2$ nor their
difference moved, so none of the three is a truncation artefact. But only the difference
has a closed form of bounded size: over the full sweep $N_1$ ranges over $[-3,5]$ and
$N_2$ over $[-2,5]$, growing with $|\bl|$ through the unbounded telescope
$\sum_{d'}A_{d'}$, while $N_1-N_2$ is confined to $[-3,3]$. I therefore write no formula
for $N_1$ alone, as the brief instructed.

\section{The closed form}\label{sec:closed}

\subsection{Cases (A) and (B): the ribbon-shift lemma}

Everything about the same-moves cases follows from one computation: what an $e$-ribbon
does to a Chevalley exponent sitting elsewhere.

\begin{lemma}[Ribbon shift]\label{lem:shift}
Let $\rho$ be a multipartition, $\gamma$ a node at $(d_0,x_0)$ addable in $\rho$, and let
the ribbon move $b\mapsto b+e$ on runner $d_1$ be legal in $\rho$; write $\tilde\rho$ for
its result. Then
\[
\Nlt(\tilde\rho,\gamma)-\Nlt(\rho,\gamma)
=\big(\delta_{x_0,b+e}-\delta_{x_0,b+e+1}\big)
+[d_1\prec d_0]\big(\delta_{x_0,b}-\delta_{x_0,b+1}-\delta_{x_0,b+e}+\delta_{x_0,b+e+1}\big).
\]
\end{lemma}

\begin{proof}
By Theorem \texttt{thm:tele},
$\Nlt(\rho,\gamma)=\sum_{d'}A_{d'}(x_0)+\sum_{d'\prec d_0}\chi_{d'}(x_0)$, and only
runner $d_1$ changes. There $\Ind_{\tilde M}-\Ind_{M}=-\delta_b+\delta_{b+e}$, so by
Lemma~\ref{lem:chi}
\[
\Delta\chi_{d_1}(y)=\big(\Ind_{\tilde M}-\Ind_M\big)(y-1)-\big(\Ind_{\tilde M}-\Ind_M\big)(y)
=\delta_{y,b}-\delta_{y,b+1}-\delta_{y,b+e}+\delta_{y,b+e+1}.
\]
The four sites lie in two residue classes: $b,b+e\equiv b$ and $b+1,b+e+1\equiv b+1$
$\pmod e$. Now
$\Delta A_{d_1}(x_0)=\sum_{k\ge1}\Delta\chi_{d_1}(x_0-ke)$, a sum over
$x_0-e,x_0-2e,\dots$. If $x_0\equiv b$, the term $+\delta_{y,b}$ contributes $+1$ exactly
when $x_0\ge b+e$ and the term $-\delta_{y,b+e}$ contributes $-1$ exactly when
$x_0\ge b+2e$; the net is $\delta_{x_0,b+e}$. If $x_0\equiv b+1$, the terms
$-\delta_{y,b+1}$ and $+\delta_{y,b+e+1}$ contribute $-1$ when $x_0\ge b+e+1$ and $+1$
when $x_0\ge b+2e+1$; the net is $-\delta_{x_0,b+e+1}$. Otherwise both classes miss and
$\Delta A_{d_1}(x_0)=0$. In all cases
$\Delta A_{d_1}(x_0)=\delta_{x_0,b+e}-\delta_{x_0,b+e+1}$. The tie term contributes
$\Delta\chi_{d_1}(x_0)$ if and only if $d_1\prec d_0$ (in particular never when
$d_1=d_0$, since $\prec$ is strict). Adding the two gives the statement.
\end{proof}

\begin{remark}
The lemma says something one could have guessed the shape of but not the support: an
$e$-ribbon is invisible to every Chevalley exponent except at the four sites
$b,b+1,b+e,b+e+1$, and the telescope sees only the \emph{upper} pair. A ribbon inserted
far below a node does not perturb its weight at all, however many beads it passes. That
is why the answer is a two-site indicator and not a count.
\end{remark}

\begin{theorem}[Case (A)]\label{thm:A}
Let $\Phi_{\nu\bl}$ be an entry whose net bead displacement moves one bead $b\mapsto b+e$
on runner $d_1$ and one bead $x_0\mapsto x_0-1$ on a runner $d_0\ne d_1$. Then it has
exactly two paths, in opposite terms, with equal ribbon heights, and
\[
N_1-N_2=\begin{cases}
\delta_{x_0,b}-\delta_{x_0,b+1}, & d_1\prec d_0,\\[2pt]
\delta_{x_0,b+e}-\delta_{x_0,b+e+1}, & d_1\succ d_0.
\end{cases}
\]
In particular $N_1-N_2\in\{-1,0,+1\}$, and the entry is nonzero exactly when $x_0$ is
one of the two indicated sites --- the lower end of the ribbon (and its successor) if the
ribbon lies on the lower runner, the upper end (and its successor) if on the upper.
\end{theorem}

\begin{proof}
The two moves are on different runners, hence independent: both orders are legal, the
pair is determined by the net displacement (a ribbon moves a bead by $+e$ and a node by
$-1$, and $e\ge2$), and there are no other paths. Put $\rho=\bl-\gamma$. In the $e_iR$
order the ribbon acts on $\bl$ and the removal has target $\nu$; in the $Re_i$ order the
removal has target $\rho$ and the ribbon then carries $\rho$ to $\nu$. Hence
$\nu=\tilde\rho$ and
\[
N_1=\Nlt(\tilde\rho,\gamma),\qquad N_2=\Nlt(\rho,\gamma).
\]
The heights are equal because the ribbon window lies on runner $d_1$ and the node move
is on $d_0\ne d_1$, so $M_{d_1}$ is the same in $\bl$ and in $\rho$. Lemma~\ref{lem:shift}
with base $\rho$ gives the difference; when $d_1\prec d_0$ the two brackets combine to
$\delta_{x_0,b}-\delta_{x_0,b+1}$, and when $d_1\succ d_0$ only the first survives. All
the data $b,d_1$ may be read in $\bl$, since $\bl$ and $\rho$ agree on runner $d_1$.
\end{proof}

\begin{theorem}[Case (B): the entry is always zero]\label{thm:B}
Let $\Phi_{\nu\bl}$ be an entry whose net displacement moves two beads on a single runner
$d$, by $b\mapsto b+e$ and $x_0\mapsto x_0-1$ with $x_0\ne b+e$ and $x_0-1\ne b$. Then
$h_1=h_2$ and $N_1=N_2$, so $\Phi_{\nu\bl}=0$. No form-(ii) entry arises this way.
\end{theorem}

\begin{proof}
Both orders are legal, and legality is more restrictive than it looks. Write
$M=M_d(\bl)$. In the $e_iR$ order the ribbon needs $b\in M$, $b+e\notin M$; then removal
at $x_0$ needs $x_0\in\tilde M$ and $x_0-1\notin\tilde M$ where
$\tilde M=M\setminus\{b\}\cup\{b+e\}$, which given $x_0\ne b+e$ and $x_0-1\ne b$ forces
$x_0\in M$, $x_0\ne b$, $x_0-1\notin M$, $x_0-1\ne b+e$. So
\[
x_0\notin\{b,\;b+1,\;b+e,\;b+e+1\},
\]
the last two by hypothesis and the first two by legality. This is precisely the support
of Lemma~\ref{lem:shift}, and since $d_1=d_0=d$ the bracket $[d_1\prec d_0]$ vanishes;
hence $N_1-N_2=\delta_{x_0,b+e}-\delta_{x_0,b+e+1}=0$.

Heights: $h_1=\#\big(M\cap(b,b+e)\big)$ and $h_2$ is the same count for
$M'=M\setminus\{x_0\}\cup\{x_0-1\}$. The two differ only if exactly one of $x_0-1,x_0$
lies in $\{b+1,\dots,b+e-1\}$, i.e.\ only if $x_0=b+1$ or $x_0=b+e$, both excluded. So
the two monomials $+t^{h_1}q^{-N_1}$ and $-t^{h_2}q^{-N_2}$ are equal and opposite.
\end{proof}

\begin{remark}
Theorem~\ref{thm:B} is the pleasant kind of result: the four sites at which a ribbon can
perturb a Chevalley exponent are exactly the four sites at which the node move would
collide with the ribbon and so are exactly the four forbidden by legality. The two moves
are far enough apart to commute precisely because they are far enough apart to both be
legal. The census (\S\ref{sec:census}) confirms the mechanism directly: over the six
configurations tabulated there the offset $x_0-b$ takes the values
$\{-6,\dots,-1\}\cup\{2,\dots,e-1\}\cup\{e+2,\dots\}$ and \emph{never} lands in
$\{0,1,e,e+1\}$.
\end{remark}

\subsection{Case (C): the $\sigma/\tau$ vocabulary, again}

\begin{theorem}[Case (C)]\label{thm:C}
Let $\Phi_{\nu\bl}$ be an entry whose net displacement moves a single bead
$c\mapsto c+e-1$ on runner $d$, with the two paths of Lemma \texttt{lem:paths} lying in
opposite terms --- equivalently, with
$\big(c-1\in M_d(\bl)\big)\Leftrightarrow\big(c+e\in M_d(\bl)\big)$. Then $h_1=h_2$ and
\[
N_1-N_2=-\varepsilon\big(\sigma(\bl;c,d)+\tau(\bl;c+e,d)\big),\qquad
\varepsilon=\begin{cases}+1,& c-1\in M_d(\bl),\\ -1,& c-1\notin M_d(\bl).\end{cases}
\]
Consequently $|N_1-N_2|\le\ell-1$, and the entry is nonzero exactly when
$\sigma(\bl;c,d)+\tau(\bl;c+e,d)\ne0$.
\end{theorem}

\begin{proof}
Write $M=M_d(\bl)$; by hypothesis $c\in M$ and $c+e-1\notin M$. By Lemma
\texttt{lem:paths} the two paths are the $c$-path (ribbon $c-1\mapsto c+e-1$, node at
content $c$) and the $(c+e)$-path (ribbon $c\mapsto c+e$, node at content $c+e$), and the
$c$-path lies in $e_iR$ iff $c-1\in M$ while the $(c+e)$-path lies in $e_iR$ iff
$c+e\notin M$. Opposite terms is therefore the stated equivalence, and $\varepsilon$ is
the sign of the case as in Theorem \texttt{thm:j}.

\emph{Heights.} With $H=\#\big(M\cap\{c+1,\dots,c+e-2\}\big)$, the height table of Lemma
\texttt{lem:paths} gives, for $\varepsilon=+1$, $h_c=H+1$ ($c$-path in $e_iR$) and
$h_{c+e}=H+1$ ($(c+e)$-path in $Re_i$); for $\varepsilon=-1$, $h_c=H$ and $h_{c+e}=H$.
Equal in both cases, so the entry is a single $t$-power.

\emph{Exponents.} Set
\[
D:=\Nlt\big(\text{base of the }(c+e)\text{-path},\ \gamma_{c+e}\big)
   -\Nlt\big(\text{base of the }c\text{-path},\ \gamma_{c}\big),
\]
each base being the target of that path's removal. Since for $\varepsilon=+1$ the
$c$-path is the $e_iR$ one and for $\varepsilon=-1$ it is the $Re_i$ one, in both cases
$N_1-N_2=-\varepsilon D$. It remains to show $D=\sigma(\bl;c,d)+\tau(\bl;c+e,d)$.

Both bases agree with $\bl$ off runner $d$. Off runner $d$, using
$A_{d'}(x+e)=\chi_{d'}(x)+A_{d'}(x)$,
\[
\sum_{d'\ne d}\big(A_{d'}(c+e)-A_{d'}(c)\big)=\sum_{d'\ne d}\chi_{d'}(c)
=\sigma(\bl;c,d)+\tau(\bl;c,d),
\]
and the two tie terms contribute $\tau(\bl;c+e,d)-\tau(\bl;c,d)$; together
$\sigma(\bl;c,d)+\tau(\bl;c+e,d)$, as required. So it suffices to show that the runner-$d$
contribution to $D$ vanishes.

\emph{Case $\varepsilon=+1$} ($c-1\in M$, $c+e\in M$). The $c$-path's base is
$\nu$ with $M_d(\nu)=M\setminus\{c\}\cup\{c+e-1\}$; the $(c+e)$-path's base is $\rho$
with $M_d(\rho)=M\setminus\{c+e\}\cup\{c+e-1\}$. The runner-$d$ contribution is
$A^{\rho}_d(c+e)-A^{\nu}_d(c)$.
\begin{itemize}
\item $A^{\nu}_d(c)$ has arguments $c-e,c-2e,\dots$, all $\le c-2$. Now $\nu$ differs
from $M$ only at $c$ and $c+e-1$, so $\chi^{\nu}_d$ differs from $\chi^{M}_d$ only at
$\{c,c+1,c+e-1,c+e\}$, all $\ge c$. Hence $A^{\nu}_d(c)=A^{M}_d(c)$.
\item $A^{\rho}_d(c+e)=\chi^{\rho}_d(c)+A^{\rho}_d(c)$. Here $\rho$ differs from $M$ only
at $c+e-1,c+e$, so $\chi^{\rho}_d$ differs from $\chi^{M}_d$ only at
$\{c+e-1,c+e,c+e+1\}$, all $>c$; hence $A^{\rho}_d(c)=A^{M}_d(c)$ and
$\chi^{\rho}_d(c)=\chi^{M}_d(c)=\Ind_M(c-1)-\Ind_M(c)=1-1=0$, using $c-1\in M$ and
$c\in M$.
\end{itemize}
So the runner-$d$ contribution is $A^{M}_d(c)-A^{M}_d(c)=0$.

\emph{Case $\varepsilon=-1$} ($c-1\notin M$, $c+e\notin M$). The $(c+e)$-path's base is
$\nu$ with $M_d(\nu)=M\setminus\{c\}\cup\{c+e-1\}$; the $c$-path's base is $\rho'$ with
$M_d(\rho')=M\setminus\{c\}\cup\{c-1\}$. The runner-$d$ contribution is
$A^{\nu}_d(c+e)-A^{\rho'}_d(c)$.
\begin{itemize}
\item $A^{\nu}_d(c+e)=\chi^{\nu}_d(c)+A^{\nu}_d(c)=\chi^{\nu}_d(c)+A^{M}_d(c)$ by the
first bullet above, and
$\chi^{\nu}_d(c)=\Ind_\nu(c-1)-\Ind_\nu(c)=0-0=0$, using $c-1\notin M$ (unchanged in
$\nu$) and $c\notin M_d(\nu)$.
\item $A^{\rho'}_d(c)$ has arguments $\le c-e\le c-2$. Now $\rho'$ differs from $M$ only
at $c-1,c$, so $\chi^{\rho'}_d$ differs from $\chi^M_d$ only at $\{c-1,c,c+1\}$, all
$\ge c-1>c-2$. Hence $A^{\rho'}_d(c)=A^{M}_d(c)$.
\end{itemize}
Again the runner-$d$ contribution is $0$. In both cases
$D=\sigma(\bl;c,d)+\tau(\bl;c+e,d)$.

The bound: $\sigma(\bl;c,d)$ is a sum of $\chi_{d'}(c)\in\{-1,0,1\}$ over the runners
above $d$ and $\tau(\bl;c+e,d)$ a sum over those below, so $|\sigma+\tau|\le\ell-1$.
\end{proof}

\begin{remark}[What the vocabulary did and did not carry]
The brief asked that the $\sigma/\tau$ vocabulary of Theorem \texttt{thm:j} be tested
before any new statistic was invented. It carries case (C) exactly --- same two half-tie
terms, same two contents $c$ and $c+e$, same bead condition $c-1\in M_d$ selecting the
sign --- and it does not carry case (A) at all. The reason is structural rather than
accidental: $\sigma$ and $\tau$ measure cross-runner traffic \emph{at the two ends of one
ribbon window}, which is a quantity only the one-runner cases have. In case (A) the node
and the ribbon are on different runners and there is no window to cross; what survives is
the bare two-site indicator of Lemma~\ref{lem:shift}. So the honest report is: the
surviving vocabulary covered half the branch, and I can say which half and why.
\end{remark}

\begin{remark}[The relation to $j$]
Comparing with Theorem \texttt{thm:j}, $j=1+\varepsilon(\sigma+\tau)$ and, in case (C),
$N_1-N_2=-\varepsilon(\sigma+\tau)$, so on the same displacement data
\[
j+(N_1-N_2)=1 .
\]
The two branches are the same computation with the two paths' terms swapped: swapping the
terms reverses the sign of the second monomial and removes the leading $1$, which is the
runner-$d$ term $\chi_d(c)=\pm1$ of Theorem \texttt{thm:j} and is exactly what the two
vanishing runner-$d$ computations above replace by $0$. Form (i) and form (ii) are not
two phenomena; they are one formula read on the two sides of a bead condition.
\end{remark}

\section{Gap 2: form (ii) is empty at level one}\label{sec:gap2}

\begin{theorem}\label{thm:l1}
At $\ell=1$ every nonzero entry of $[e_i,R^{(1)}(t)]$ is of form (i), with $j=1$. Form
(ii) is empty.
\end{theorem}

\begin{proof}
A path consists of one ribbon and one node move, so the net displacement $\Delta$ moves
at most two beads, and the cases of Theorem \texttt{thm:dich} are exhaustive: $\Delta$
moves beads on two runners (case (A)), or two beads on one runner (case (B)), or one bead
on one runner (form (i) or case (C)). ($\Delta=0$ is impossible: a ribbon displaces a bead
by $+e$ and a node by $-1$, and these cancel only if $e=1$.)

At $\ell=1$: case (A) cannot occur, there being one runner. Case (B) gives $N_1=N_2$ by
Theorem~\ref{thm:B}, at any level, so the entry is zero. Case (C) has
$\sigma(\bl;c,d)=\tau(\bl;c+e,d)=0$, both being empty sums over the runners other than
$d$, so $N_1-N_2=0$ by Theorem~\ref{thm:C} and the entry is again zero. Hence every
nonzero entry has its two paths in the same term, which is form (i), and there
$j=1+\varepsilon\cdot0=1$.
\end{proof}

\begin{remark}[Not a vacuity]
\emph{Variation exercised:} at $\ell=1$, do case-(B) and case-(C) path pairs occur at
all? If they did not, Theorem~\ref{thm:l1} would be true for a reason having nothing to
do with the cancellation it claims. They occur in quantity:
\begin{center}
\begin{tabular}{lrrrr}
\hline
configuration ($\ell=1$, $|\lambda|\le7$) & case-(B) pairs & case-(C) pairs & total & nonzero\\
\hline
$e=2$, $s=0$ & $142$ & $76$ & $218$ & $0$\\
$e=3$, $s=0$ & $182$ & $88$ & $270$ & $0$\\
$e=4$, $s=0$ & $224$ & $104$ & $328$ & $0$\\
$e=5$, $s=0$ & $278$ & $112$ & $390$ & $0$\\
$e=2$, $s=1$ & $142$ & $76$ & $218$ & $0$\\
\hline
total & $968$ & $456$ & $1424$ & $0$\\
\hline
\end{tabular}
\end{center}
All $1424$ opposite-term pairs at $\ell=1$ have $N_1=N_2$ and $h_1=h_2$, and every one of
them cancels. The emptiness is a cancellation over a live population, which is what the
theorem asserts.
\end{remark}

\begin{remark}[Independent confirmation, with a positive control]
Separately, the full classification was rerun at $\ell=1$ at sizes beyond every previous
sweep --- $(e,s,|\lambda|_{\max})=(2,0,9),(3,0,8),(4,0,8),(5,0,8),(2,3,8),(3,1,8)$, bead
depths $18$--$24$. It returns $607$ form-(i) entries, $0$ form-(ii) entries and $0$
unclassified entries, with $j=1$ on all $607$. The $607$ are the positive control: a run
that found nothing at all would have proved nothing about form (ii).
\end{remark}

\section{The classification, complete}\label{sec:class}

\begin{theorem}[Classification of $[e_i,R^{(\ell)}(t)]$]\label{thm:full}
Let $\Phi_{\nu\bl}(t)\ne0$ be an entry of $[e_i,R^{(\ell)}(t)]$ and $\Delta$ the net
change of Maya sets. Then $\Phi_{\nu\bl}$ has exactly two contributing paths and exactly
one of the following holds.
\begin{enumerate}
\item[(i)] $\Delta$ moves one bead $c\mapsto c+e-1$ on a runner $d$, and
$c-1\in M_d(\bl)\Leftrightarrow c+e\notin M_d(\bl)$. Then the two paths lie in the same
term and
$\Phi=\varepsilon q^{k}t^{b}\big(1+q^{j}t\big)$ with
$j=1+\varepsilon\big(\sigma(\bl;c,d)+\tau(\bl;c+e,d)\big)$ and $\varepsilon=+1$ iff
$c-1\in M_d(\bl)$ \emph{(Theorem \texttt{thm:j})}.
\item[(ii-C)] $\Delta$ moves one bead $c\mapsto c+e-1$ on a runner $d$, and
$c-1\in M_d(\bl)\Leftrightarrow c+e\in M_d(\bl)$. Then the two paths lie in opposite
terms and $\Phi=q^{-N_2}t^{h}\big(q^{-m}-1\big)$ with
$m=-\varepsilon\big(\sigma(\bl;c,d)+\tau(\bl;c+e,d)\big)\ne0$; here $1\le|m|\le\ell-1$.
\item[(ii-A)] $\Delta$ moves one bead $b\mapsto b+e$ on a runner $d_1$ and one bead
$x_0\mapsto x_0-1$ on a runner $d_0\ne d_1$, with $x_0=b$ or $b+1$ if $d_1\prec d_0$, and
$x_0=b+e$ or $b+e+1$ if $d_1\succ d_0$. Then
$\Phi=q^{-N_2}t^{h}\big(q^{\mp1}-1\big)$, the sign being $-$ at the first of the two
sites and $+$ at the second.
\end{enumerate}
No other $\Delta$ gives a nonzero entry: two beads on one runner always cancel
\emph{(Theorem~\ref{thm:B})}, as does case (ii-C) with $\sigma+\tau=0$ and case (ii-A)
with $x_0$ at neither site. In particular \emph{(Theorem~\ref{thm:l1})} at $\ell=1$ only
(i) survives, with $j=1$.
\end{theorem}

\begin{proof}
Immediate from Theorem \texttt{thm:dich}, Theorem \texttt{thm:j}, and
Theorems~\ref{thm:A}, \ref{thm:B}, \ref{thm:C} above.
\end{proof}

\begin{corollary}\label{cor:q1}
At $q=1$ every form-(ii) entry vanishes and every form-(i) entry becomes
$2\varepsilon t^{b}$.
\end{corollary}

\begin{corollary}[Q59 recovered]\label{cor:q59}
At $\ell=1$ every nonzero entry of $[e_i,R^{(1)}(t)]$ equals $\varepsilon q^{k}t^{b}(1+qt)$.
\end{corollary}
\begin{proof}
By Theorem~\ref{thm:l1} every nonzero entry is of form (i) with $j=1$, and form (i) with
$j=1$ is $\varepsilon q^{k}t^{b}(1+qt)$.
\end{proof}

This is Corollary \texttt{cor:Q59} of 31 August, whose content is that $(1+qt)$ divides
$[e_i,R_e(t)]$ entrywise, hence as a matrix: $[e_i,R_e(t)]=(1+qt)D^{-}_i(t)$. So the
level-$\ell$ classification degenerates to the level-$1$ rigidity theorem, as it must,
and the divisibility fails at higher level for a reason the classification now names ---
$j\ne1$ in case (i), and the form-(ii) entries, which are $t$-monomials and are divisible
by $(1+qt)$ only when they are zero.

\section{Verification and negative controls}\label{sec:controls}

\paragraph{Same object.} The nine in-sample triples return $1248$ form-(i) entries and
$0$ unclassified entries, reproducing 2 September exactly; the $1604$ form-(ii) entries
are the complement.

\paragraph{In sample.} Theorems~\ref{thm:A} and \ref{thm:C}: $\mathbf{1604/1604}$.

\paragraph{Out of sample.} Seven configurations absent from the in-sample list ---
$(e,\ell,\bs)=(2,2,(1,0))$, $(3,2,(0,2))$, $(2,4,(0,0,0,0))$, $(5,2,(0,0))$,
$(3,3,(0,1,2))$, $(2,2,(0,0))$ at $|\bl|\le5$, and $(2,3,(0,0,0))$ at depth $+6$ ---
giving $1564$ further form-(ii) entries. $\mathbf{1564/1564}$. The $\ell=4$ configuration
produces $|N_1-N_2|=3$, a value no in-sample run had seen, predicted correctly; this is
the form-(ii) analogue of the $j=4$ event of 2 September.

\paragraph{Depth.} The standing instruction is to check depth first on any failure.
There were no failures, so the control was run positively: $(2,3,(0,0,0))$ recomputed
with the bead depth raised by $6$ returns the same $354$ form-(ii) entries with the same
$N_1$, $N_2$ and $N_1-N_2$ throughout, and the separate depth-sensitivity run of
\S\ref{sec:census} moved none of $590$ tracked $(N_1,N_2)$ pairs. The counts are not
truncation artefacts.

\paragraph{Range.} $|N_1-N_2|\le1$ in case (A) at every level tested, and
$|N_1-N_2|\le\ell-1$ in case (C), \emph{saturated} at each of $\ell=2,3,4$.

\paragraph{Sign identification.} All $3168$ entries have coefficient exactly
$q^{-N_1}-q^{-N_2}$ with $N_1$ from the $e_iR$ path, confirming the labelling on which
the sign of the theorem depends.

\paragraph{Negative controls.} Each was first checked for non-degeneracy --- it must
change at least one prediction --- and the count of moved predictions is reported
beside its score. \emph{Every variant below moves a nonzero number of predictions;
none is blind.} Scored on all $3168$ entries.

\begin{center}
\begin{tabular}{lrrl}
\hline
variant & score & moved & variation exercised\\
\hline
Theorems~\ref{thm:A}, \ref{thm:C} & $\mathbf{3168/3168}$ & --- & ---\\
(A) the two runner-order branches exchanged & $872$ & $2296$ & the order of $d_1$ vs $d_0$\\
(A) runner order ignored (always the $b+e$ branch) & $2020$ & $1148$ & whether $d_1$ vs $d_0$ matters\\
(A) the two-site sign flipped & $872$ & $2296$ & which site carries $+1$\\
(C) $\varepsilon$ branch ignored & $2732$ & $436$ & the bead condition $c-1\in M_d$\\
(C) $\sigma\leftrightarrow\tau$ exchanged & $2434$ & $734$ & tie-break orientation\\
(C) both halves read at content $c$ & $2738$ & $430$ & which content each half uses\\
(C) both halves read at content $c+e$ & $2738$ & $430$ & which content each half uses\\
(C) the $\sigma$ term dropped & $2717$ & $451$ & the runners above $d$\\
(C) the $\tau$ term dropped & $2717$ & $451$ & the runners below $d$\\
(C) overall sign flipped & $2296$ & $872$ & orientation of $e_iR-Re_i$\\
cross-case: (A)'s rule applied to (C) too & $2296$ & $872$ & the case split itself\\
cross-case: (C)'s rule applied to (A) too & $872$ & $2296$ & the case split itself\\
rival: sign read off the $c$-path's term alone & $1366$ & $1802$ & a cheaper statistic\\
\textbf{floor:} the constant $N_1-N_2\equiv1$ & $1530$ & $1638$ & --- \\
\hline
\end{tabular}
\end{center}

The floor is $1530/3168$ ($48.3\%$), which is what a constant can score on a
sign-symmetric quantity. Every structural ingredient of the theorem is worth between
$430$ and $2296$ moved predictions, and every perturbation of one loses.

\begin{remark}[A control I wrote twice by mistake, and the correction]
My first control list contained ``$\sigma\leftrightarrow\tau$ exchanged'' and ``both
halves read at the other node's content'' as two entries. They are the same expression:
sending $\sigma(c)+\tau(c+e)\mapsto\tau(c)+\sigma(c+e)$ \emph{is} both operations at once.
The table above replaces the second by two genuinely different variants (both halves at
$c$, both at $c+e$). This is the same failure mode as the blind control of 2 September in
a smaller key --- a control that is not the control you named --- and it is recorded
rather than quietly fixed.
\end{remark}

\paragraph{Planted errors.} Perturbing $1$, $5$ and $20$ of the \emph{observed}
$N_1-N_2$ by $+1$ drops the score to $3167$, $3163$ and $3148$: caught one for one, so
the comparison can see a failure.

\paragraph{What else would pass.} There is no fitted parameter. In case (A) the four
sites are the support computed in Lemma~\ref{lem:shift} and the branch is the bracket
$[d_1\prec d_0]$ appearing in Theorem \texttt{thm:tele}; in case (C) $\sigma$ and $\tau$
are the two halves of that same tie term, the contents $c$ and $c+e$ are forced by Lemma
\texttt{lem:paths}, and $\varepsilon$ is the bead condition already proved to give the
sign in Theorem \texttt{thm:j}. The one place a coincidence could hide is the case (A)
branch, whose two halves are distinguished only on the $1148$ entries with $d_1\prec d_0$
--- which is why the ``runner order ignored'' control is reported: it still scores
$2020$, and the branch is supported by the $1148$ entries it moves.

\section{Gaps}\label{sec:gaps}

\begin{enumerate}
\item \textbf{Sizes.} $|\bl|\le5$ at $\ell\le3$, $\le3$ at $\ell=4$; bead depths
$12$--$24$. $\ell\le4$, so $\sigma$ and $\tau$ have been exercised with at most three
other runners. The $\ell=1$ runs reach $|\lambda|\le9$.

\item \textbf{One tie-break.} As on 2 September, everything is computed at tie-break
$+1$. Under the opposite convention $\sigma$ and $\tau$ exchange roles and so do the two
branches of Theorem~\ref{thm:A}; the mirror statement is not written out. The
$\sigma\leftrightarrow\tau$ control measures the difference and it is large ($734$
predictions), so the convention is doing real work and the mirror is not free.

\item \textbf{$N_1$ and $N_2$ separately.} Only the difference is given a closed form.
$N_1$ alone is $\sum_{d'}A_{d'}+\tau$, an unbounded telescope; no shorter description of
it is claimed or attempted. See \S\ref{sec:census}, census question 3.

\item \textbf{Proposition \texttt{prop:heis} stands, untouched.} As stated in the
abstract: for $\ell\ge2$, $R^{(\ell)}(-q^{-1})\ne B^{[e]}_{-1}$, so
Theorem~\ref{thm:full} classifies the commutator of a concretely defined operator and is
not known to be a statement about the level-$\ell$ Heisenberg action. This is the
programme's largest open question and no session since 31 August has moved it.

\item \textbf{No representation-theoretic reading of the two form-(ii) cases.} The
theorem says two-runner entries are always $\pm1$ and one-runner entries carry the
cross-runner traffic, but I have no conceptual account of why the two-runner case is
rigid. Lemma~\ref{lem:shift} explains it mechanically --- a ribbon is invisible except at
four sites --- without saying what it means.
\end{enumerate}

\section{Corrections to earlier papers}\label{sec:corr}

\begin{enumerate}
\item Gap 1 of 2 September (``form (ii) is not computed'') is \textbf{closed}
(Theorems~\ref{thm:A}, \ref{thm:B}, \ref{thm:C}).
\item Gap 2 of 2 September (``form (ii) is not proved empty at $\ell=1$'') is
\textbf{closed} (Theorem~\ref{thm:l1}); the Corollary at the end of
\S\texttt{sec:dichotomy}, whose $\ell=1$ emptiness clause was explicitly left
computational, is now proved in full.
\item The 2 September speculation that ``the same $\sigma/\tau$ analysis should apply''
to form (ii) is \textbf{half right}: it applies verbatim to case (C) and not at all to
case (A), which needs Lemma~\ref{lem:shift} instead. Recorded because the speculation was
written as though the branch were homogeneous, and it is not.
\end{enumerate}

\section*{Is this a paper?}

Yes, with one named weakness. What now exists is a complete and self-contained
classification of $[e_i,R^{(\ell)}(t)]$ for the graded ribbon operator
$R^{(\ell)}(t)$ on level-$\ell$ Fock space: every nonzero entry is one of three explicit
shapes, each shape's exponent is given by a formula in the multipartition and the net
bead displacement alone, every formula is proved from the telescope theorem rather than
fitted, and the level-one case reduces to the Q59 rigidity statement as it must. It rests
on three theorems proved across three sessions (\texttt{thm:tele}, \texttt{thm:j},
Theorem~\ref{thm:full}) with no remaining computational premises, it has $3168$
out-of-sample-inclusive confirmations and a floor to measure them against, and its
negative controls are non-degenerate by construction. A referee could read it.

The single weakest point is not in the mathematics: it is the motivation. Proposition
\texttt{prop:heis} says that for $\ell\ge2$ the operator being classified is not the one
the programme wanted --- $R^{(\ell)}(-q^{-1})$ is not $B^{[e]}_{-1}$ --- so a reader who
does not already care about $R^{(\ell)}(t)$ has, at present, no reason given here to
start. At $\ell=1$ the operator \emph{is} the right one and the paper's level-one content
is a genuine statement about the Heisenberg action; at higher level the paper is
currently a theorem about a well-defined combinatorial operator whose representation
theory is unestablished. Anyone deciding whether to promote this to
\texttt{publishable-result} should decide that question first, because it determines
whether the paper is ``the commutator of the level-$\ell$ ribbon operator, classified''
or ``a level-$\ell$ analogue of Q59, with the analogy unverified''. The honest framing is
the second, and the second is still worth publishing --- but it must be the title's
framing, not a remark's.

\end{document}
